\documentclass[a4paper,11pt]{article}

\usepackage[T1]{fontenc}
\usepackage[utf8]{inputenc}
\usepackage[top=3cm,bottom=3cm,left=2.5cm,right=2.5cm]{geometry}
\usepackage{lmodern}
\usepackage{microtype}
\usepackage{amsmath}
\usepackage{booktabs}
\usepackage{longtable}
\usepackage{array}
\usepackage{xcolor}
\usepackage{fancyhdr}
\usepackage[hidelinks,
            pdftitle={CPM-IA Run Report},
            pdfauthor={Math Biology S.r.l.},
            pdfsubject={SR\&TS \#L001-U015-P26.0127}]{hyperref}

%% ─── Colours ────────────────────────────────────────────────────────────────
\definecolor{dmaaccent}{RGB}{0,74,128}
\definecolor{dmagray}{RGB}{100,100,100}

%% ─── Page style ─────────────────────────────────────────────────────────────
\pagestyle{fancy}
\fancyhf{}
\lhead{\small\textcolor{dmagray}{Math Biology S.r.l.\ --- Confidential}}
\rhead{\small\textcolor{dmagray}{CPM-IA Automated Run Report}}
\lfoot{\small\textcolor{dmagray}{SR\&TS \#L001-U015-P26.0127 --- Rev~00.00}}
\rfoot{\small\textcolor{dmagray}{Page~\thepage}}
\renewcommand{\headrulewidth}{0.4pt}
\renewcommand{\footrulewidth}{0.4pt}

%% ─── Section colour ──────────────────────────────────────────────────────────
\usepackage{titlesec}
\titleformat{\section}{\large\bfseries\color{dmaaccent}}{}{0em}{}[\titlerule]
\titleformat{\subsection}{\normalsize\bfseries\color{dmaaccent}}{}{0em}{}

%% ─────────────────────────────────────────────────────────────────────────────
\begin{document}

%% ── Title page ───────────────────────────────────────────────────────────────
\begin{titlepage}
\begin{center}
  \vspace*{1.8cm}
  {\color{dmaaccent}\rule{\linewidth}{1.5pt}}\\[0.6cm]
  {\LARGE\bfseries\color{dmaaccent}
    Cross Points and Markers\\[0.25cm]
    Influence Analysis}\\[0.4cm]
  {\Large Automated Run Report}\\[0.3cm]
  {\color{dmaaccent}\rule{\linewidth}{1.5pt}}\\[2cm]

  \begin{tabular}{ll}
    \textbf{Component ID:}   & CPM-IA \\[0.3cm]
    \textbf{SR\&TS Ref.:}    & \#L001-U015-P26.0127 --- Rev~00.00 \\[0.3cm]
    \textbf{Run date:}       & 2026-08-31 \\[0.3cm]
    \textbf{Run time (UTC):} & 13:48:49 \\
  \end{tabular}

  \vfill
  {\color{dmagray}\small
    Math Biology S.r.l.\ --- All rights reserved.\\
    This document is generated automatically at pipeline completion.\\
    Do not distribute without authorisation.}
\end{center}
\end{titlepage}

\tableofcontents
\clearpage

%% ── 1. Run Configuration ────────────────────────────────────────────────────
\section{Run Configuration}

\begin{table}[h!]
\centering
\small
\caption{Runtime parameters loaded from the XML configuration file.}
\begin{tabular}{ll}
\toprule
\textbf{Parameter} & \textbf{Value} \\
\midrule
Input path                        & \texttt{data/input/all\_visits\_percentage\_variation\_cleaned.csv} \\
Output path                       & \texttt{data/output} \\
Detail output file                & \texttt{cpm\_ia\_detail.csv} \\
Aggregated output file            & \texttt{cpm\_ia\_aggregated.csv} \\
Threshold $\theta$ (\%)           & 300.0 \\
Rise tolerance $\varepsilon$ (\%) & 5.0 \\
Max descent window $n_{\max}$     & 10 \\
\bottomrule
\end{tabular}
\end{table}

%% ── 2. Dataset Overview ─────────────────────────────────────────────────────
\section{Dataset Overview}

\begin{table}[h!]
\centering
\small
\caption{Input dataset dimensions after schema validation and ingestion (Module~02).}
\begin{tabular}{lr}
\toprule
\textbf{Metric} & \textbf{Count} \\
\midrule
Unique visits                          & 8577 \\
Unique anatomical points               & 62 \\
Unique markers across all visits       & 2237 \\
(Visit,\,point) series analysed        & 84042 \\
Positive (visit,\,point) series        & 15912 (\,18.9\,\%) \\
Descent windows truncated at $n_{\max}$ & 0 (\,0.0\,\%) \\
\bottomrule
\end{tabular}
\end{table}

%% ── 3. Algorithm Description ─────────────────────────────────────────────────
\section{Algorithm Description}

The CPM-IA component analyses whether excitation events detected in bioelectrical surface
current signals systematically influence the signal values recorded at subsequent measurement
markers. The pipeline processes long-format tabular measurements through eight sequential
analytical modules (M03--M10), each producing typed, immutable outputs consumed by the next
stage.

\subsection{Input Data}

The component ingests a single long-format CSV file (UTF-8, comma-separated). Each row
encodes one (visit, anatomical point, marker) measurement triplet. The required columns are:

\begin{itemize}
  \item \textbf{Visit identifier} --- unique alphanumeric code for the acquisition session.
  \item \textbf{Marker identifier} --- unique code identifying the stimulus or measurement
        marker.
  \item \textbf{Anatomical point identifier} --- code identifying the electrode or
        body-surface acquisition site.
  \item \textbf{Percentage variation} --- raw signal change relative to a reference value,
        expressed as a floating-point percentage (\%).
\end{itemize}

Additional visit-level metadata columns (e.g.\ acquisition date) are preserved in the
detail output without modification.

\subsection{Configuration Parameters}

Three parameters govern the analysis; all are loaded exclusively from the XML configuration
file at run time and are never hard-coded:

\begin{itemize}
  \item \textbf{Threshold} $\theta = \mathbf{300.0}\,\%$: a marker is
        classified as a \emph{positive excitation} when its percentage variation strictly
        exceeds $\theta$.
  \item \textbf{Rise tolerance} $\varepsilon = \mathbf{5.0}\,\%$: the maximum
        increase permitted between two consecutive values inside the descent window; a rise
        strictly exceeding $\varepsilon$ closes the window.
  \item \textbf{Maximum window length} $n_{\max} = \mathbf{10}$: the descent window
        is forcibly closed after $n_{\max}$ markers even when no rise exceeding $\varepsilon$
        has occurred (truncated window).
\end{itemize}

\subsection{Processing Steps}

\paragraph{M03 --- Marker Sequence Reconstruction.}
For each visit independently, the canonical marker detection order is derived from the
first-appearance order of marker identifiers in the input file's rows for that visit.
This per-visit temporal ordering is immutable for the entire run: no sorting, grouping, or
reordering of markers is performed after this step.
Each (visit, anatomical~point) pair is represented as an ordered list of percentage
variation values aligned to its visit's detection order; positions where a marker belongs
to the visit protocol but was not measured at a specific point receive a NaN sentinel value.

\paragraph{M04 --- First-Positive Trigger Detection.}
Each (visit, point) sequence is scanned in temporal detection order.
The first marker whose percentage variation \emph{strictly exceeds} $\theta$ is designated
the \emph{excitation anchor} --- the onset of the excitation event.
The anchor is immutable for the run: once assigned, it cannot be reassigned even if a later
marker has a higher value.
Series in which no marker exceeds $\theta$ are flagged as ``no cross-marker effect'' and
are excluded from all subsequent analytical modules (M06--M09).

\paragraph{M05 --- Positive Census.}
An anatomical point is classified as \emph{positive} if at least one of its visits produced
an anchor. The census counts positive and total unique anatomical points across the dataset.

\paragraph{M06 --- Descent Window Determination.}
Starting from the anchor (step $i = 0$), the window extends to successive markers as long
as the signal does not rise by more than $\varepsilon$ between consecutive steps.
Formally, the window closes at the first step $i \geq 1$ where
$v_i - v_{i-1} > \varepsilon$.
The window is also closed when it reaches $n_{\max}$ markers (truncated window).
The anchor is always included as the first element of the window.

\paragraph{M07 --- Descent Slope (OLS Linear Regression).}
An ordinary least-squares (OLS) linear regression is fitted to the $L$ descent window
values $y_0, y_1, \ldots, y_{L-1}$ at integer step positions $x_i = i$, where
$x_0 = 0$ corresponds to the anchor.
The closed-form slope estimator is:
\[
  \hat{\beta}_1 =
  \frac{L\sum_{i=0}^{L-1} i\,y_i
        - \sum_{i=0}^{L-1} i \cdot \sum_{i=0}^{L-1} y_i}
       {L\sum_{i=0}^{L-1} i^2
        - \Bigl(\sum_{i=0}^{L-1} i\Bigr)^{\!2}}
\]
The slope $\hat{\beta}_1$ (in \%\,step$^{-1}$) quantifies the average linear rate of
change of the signal across the descent window. A window of length $L = 1$ yields an
undefined slope, reported as missing.

\paragraph{M08 --- Descent Descriptors.}
Three scalar shape descriptors are computed for each positive (visit, point) series:
\begin{itemize}
  \item \textbf{Descent depth} $= v_{\mathrm{anchor}} - \min(v_0, \ldots, v_{L-1})$,
        in \%: total signal drop from the anchor to the lowest point in the window.
  \item \textbf{Descent length} $= L$: number of markers in the window, including the
        anchor.
  \item \textbf{Mean per-step decrease} $= \mathrm{depth}\,/\,(L - 1)$,
        in \%\,step$^{-1}$: average drop per marker step, normalising depth by window
        duration. Undefined for $L = 1$.
\end{itemize}

\paragraph{M09 --- Variance Indicator.}
Sample variance (denominator $n - 1$, i.e.\ $\mathtt{ddof} = 1$) is computed separately
on two signal segments:
\begin{itemize}
  \item \textbf{Before-segment}: all sequence values at positions before the anchor index,
        after removing NaN entries.
  \item \textbf{After-segment}: the descent window values $v_0, \ldots, v_{L-1}$, after
        removing NaN entries.
\end{itemize}
A segment with fewer than two clean values yields an undefined variance (reported as
missing). The \textbf{variance ratio}
$= \sigma^2_{\mathrm{after}} / \sigma^2_{\mathrm{before}}$
is computed when both variances are defined and $\sigma^2_{\mathrm{before}} > 0$.
A ratio $> 1$ indicates that signal variability \emph{increases} after the excitation
onset; a ratio $< 1$ indicates it \emph{decreases}.

\paragraph{M10 --- Output Consolidation.}
Results from M04--M09 are assembled into two tidy output files:
\begin{itemize}
  \item \textbf{Detail file} (\texttt{cpm\_ia\_detail.csv}): one row per
        (visit, anatomical~point) pair, containing all analytical outputs alongside visit
        metadata.
  \item \textbf{Aggregated file} (\texttt{cpm\_ia\_aggregated.csv}): one row per
        anatomical point, containing visit counts, positive prevalence, and median descent
        descriptors computed over positive visits only.
\end{itemize}

\subsection{Output Artefacts}

The pipeline produces three artefacts in the output directory \texttt{data/output}:
\begin{enumerate}
  \item The per-(visit, point) detail CSV (M10).
  \item The per-point aggregation CSV (M10).
  \item This automated PDF run report (M11), generated from a filled \LaTeX{} template
        and compiled with \texttt{pdflatex} in two passes to resolve table-of-contents
        page numbers.
\end{enumerate}
No output file is ever silently overwritten: the pipeline raises an error if a target
file already exists (guard~G-10).

%% ── 4. Per-Point Results ────────────────────────────────────────────────────
\section{Per-Point Results}

\noindent
Of the \textbf{84042} (visit,\,point) series analysed,
\textbf{15912} (\textbf{18.9}\,\%) yielded at least one
excitation event above the threshold $\theta = 300.0\,\%$.
\textbf{62} of \textbf{62} unique anatomical points
have at least one positive visit.
All tables in this section are sorted alphabetically by point identifier.

\subsection{Visit Counts and Positive Prevalence}

A visit is \emph{positive} for a given anatomical point when at least one marker in that
(visit, point) series has a percentage variation strictly exceeding
$\theta = 300.0\,\%$ in temporal detection order (Module~M04).
All other visits are classified as ``no cross-marker effect'' for that point.

The \textbf{positive prevalence} is the proportion of total visits that are positive:
\[
  \text{Prevalence (\%)} = \frac{\text{Positive Visits}}{\text{Total Visits}} \times 100
\]

A high prevalence indicates that the excitation event is consistently reproduced at a given
point across visits; a low prevalence indicates that the event is rare or confined to
specific sessions.

\begin{small}
\begin{longtable}{p{5cm}rrr}
\toprule
\textbf{Anatomical Point} &
\textbf{Total Visits} &
\textbf{Positive Visits} &
\textbf{Prevalence (\%)} \\
\midrule
\endfirsthead
\toprule
\textbf{Anatomical Point} &
\textbf{Total Visits} &
\textbf{Positive Visits} &
\textbf{Prevalence (\%)} \\
\midrule
\endhead
\midrule
\multicolumn{4}{r}{\small\emph{Continued on next page\ldots}} \\
\endfoot
\bottomrule
\endlastfoot
LF - 1 - c [43] & 902 & 73 & 8.1 \\
LF - 1 - e [32] & 1801 & 365 & 20.3 \\
LF - 1 - i [31] & 1205 & 210 & 17.4 \\
LF - 2 - c [44] & 1137 & 178 & 15.7 \\
LF - 2 - e [34] & 1547 & 254 & 16.4 \\
LF - 2 - i [33] & 994 & 103 & 10.4 \\
LF - 3 - c [36] & 1353 & 216 & 16.0 \\
LF - 3 - e [37] & 1414 & 209 & 14.8 \\
LF - 3 - i [35] & 1585 & 313 & 19.7 \\
LF - 4 - c [39] & 1255 & 169 & 13.5 \\
LF - 4 - e [40] & 1206 & 131 & 10.9 \\
LF - 4 - i [38] & 1324 & 192 & 14.5 \\
LF - 5 - c [45] & 1359 & 115 & 8.5 \\
LF - 5 - e [42] & 1305 & 144 & 11.0 \\
LF - 5 - i [41] & 1205 & 193 & 16.0 \\
LH - 1 - c [2] & 1602 & 386 & 24.1 \\
LH - 1 - e [3] & 1522 & 376 & 24.7 \\
LH - 1 - i [1] & 1055 & 159 & 15.1 \\
LH - 2 - c [5] & 1907 & 588 & 30.8 \\
LH - 2 - e [6] & 1365 & 283 & 20.7 \\
LH - 2 - i [4] & 1942 & 534 & 27.5 \\
LH - 3 - c [8] & 1249 & 263 & 21.1 \\
LH - 3 - e [9] & 1864 & 528 & 28.3 \\
LH - 3 - i [7] & 1195 & 241 & 20.2 \\
LH - 4 - c [11] & 1464 & 375 & 25.6 \\
LH - 4 - e [12] & 1261 & 242 & 19.2 \\
LH - 4 - i [10] & 1560 & 446 & 28.6 \\
LH - 5 - c [15] & 1477 & 375 & 25.4 \\
LH - 5 - e [14] & 2155 & 650 & 30.2 \\
LH - 5 - i [13] & 1480 & 349 & 23.6 \\
RF - 1 - c [58] & 913 & 82 & 9.0 \\
RF - 1 - e [47] & 1555 & 268 & 17.2 \\
RF - 1 - i [46] & 1442 & 179 & 12.4 \\
RF - 2 - c [59] & 990 & 67 & 6.8 \\
RF - 2 - e [49] & 1421 & 205 & 14.4 \\
RF - 2 - i [48] & 969 & 76 & 7.8 \\
RF - 3 - c [51] & 1295 & 201 & 15.5 \\
RF - 3 - e [52] & 1332 & 183 & 13.7 \\
RF - 3 - i [50] & 1545 & 257 & 16.6 \\
RF - 4 - c [54] & 1267 & 136 & 10.7 \\
RF - 4 - e [55] & 1167 & 133 & 11.4 \\
RF - 4 - i [53] & 1421 & 197 & 13.9 \\
RF - 5 - c [60] & 1248 & 125 & 10.0 \\
RF - 5 - e [57] & 1362 & 179 & 13.1 \\
RF - 5 - i [56] & 1140 & 119 & 10.4 \\
RH - 1 - c [17] & 1625 & 375 & 23.1 \\
RH - 1 - e [18] & 1371 & 317 & 23.1 \\
RH - 1 - i [16] & 1124 & 154 & 13.7 \\
RH - 2 - c [20] & 1854 & 514 & 27.7 \\
RH - 2 - e [21] & 1164 & 190 & 16.3 \\
RH - 2 - i [19] & 1870 & 457 & 24.4 \\
RH - 3 - c [23] & 1135 & 187 & 16.5 \\
RH - 3 - e [24] & 1463 & 294 & 20.1 \\
RH - 3 - i [22] & 1182 & 193 & 16.3 \\
RH - 4 - c [26] & 1397 & 360 & 25.8 \\
RH - 4 - e [27] & 1207 & 218 & 18.1 \\
RH - 4 - i [25] & 1518 & 391 & 25.8 \\
RH - 5 - c [30] & 1359 & 317 & 23.3 \\
RH - 5 - e [29] & 1924 & 519 & 27.0 \\
RH - 5 - i [28] & 1645 & 327 & 19.9 \\
Thyroid - LH [61] & 237 & 22 & 9.3 \\
Thyroid - RH [62] & 210 & 10 & 4.8 \\
\end{longtable}
\end{small}

\subsection{Most Frequent Trigger Markers by Anatomical Point}

For each anatomical point, the table shows which markers most frequently acted as the
\emph{excitation anchor} --- the first marker in temporal detection order (Module~M03)
whose percentage variation strictly exceeded $\theta = 300.0\,\%$
(Module~M04) --- across positive visits.

\medskip
Only markers that acted as anchor in at least \textbf{3.5\,\%}
of the positive visits for that point are listed.
Markers with a lower share are scattered across sessions with no single dominant trigger
and are omitted for readability.
All listed marker values exceeded $\theta = 300.0\,\% > 100\,\%$ in the
visits where they acted as anchor.
Anatomical points whose most frequent anchor marker appears in fewer than
3.5\,\% of positive visits do not appear in this table (no single
marker is dominant at those points).

\medskip
\noindent\textbf{How to read the columns:}
\begin{itemize}
  \item \textbf{Trigger Marker}: identifier of the first marker to exceed $\theta$ for
        that point; the value that classified the visit as positive.
  \item \textbf{$n$}: number of positive visits for that point in which this marker was
        the anchor.
  \item \textbf{Share of Pos.\ Visits (\%)}: $n$ expressed as a percentage of the total
        positive visits for that point.
        Values near 100\,\% indicate a dominant, reproducible trigger;
        lower values indicate concurrent triggers across sessions (e.g.\ due to
        protocol or biological variability).
\end{itemize}

\medskip
\noindent Rows are sorted by anatomical point (alphabetical) and, within each point, by
descending $n$.

\begin{footnotesize}
\begin{longtable}{p{3.5cm}p{7.5cm}rr}
\toprule
\textbf{Anatomical Point} &
\textbf{Trigger Marker} &
\textbf{$n$} &
\textbf{Share (\%)} \\
\midrule
\endfirsthead
\toprule
\textbf{Anatomical Point} &
\textbf{Trigger Marker} &
\textbf{$n$} &
\textbf{Share (\%)} \\
\midrule
\endhead
\midrule
\multicolumn{4}{r}{\footnotesize\emph{Continued on next page\ldots}} \\
\endfoot
\bottomrule
\endlastfoot
LF - 1 - c [43] & Acetilcolina [1166] & 8 & 11.0 \\
LF - 1 - c [43] & Calcium Fluoratum [1151] & 8 & 11.0 \\
LF - 1 - c [43] & Bombesina [1076] & 5 & 6.8 \\
LF - 1 - c [43] & C-RAS [1163] & 5 & 6.8 \\
LF - 1 - c [43] & 3-Nitro Tirosina [361] & 3 & 4.1 \\
LF - 1 - c [43] & Insulina [2282] & 3 & 4.1 \\
LF - 1 - c [43] & Profito CIRC [639] & 3 & 4.1 \\
LF - 1 - c [43] & S-100 [1123] & 3 & 4.1 \\
LF - 1 - c [43] & TXA - 2 [551] & 3 & 4.1 \\
LF - 1 - e [32] & Candida albicans [556] & 27 & 7.4 \\
LF - 1 - e [32] & COX-2 [679] & 14 & 3.8 \\
LF - 1 - e [32] & IFN-alfa [2258] & 13 & 3.6 \\
LF - 1 - i [31] & ACTH [893] & 16 & 7.6 \\
LF - 1 - i [31] & 8-OHdG [526] & 11 & 5.2 \\
LF - 1 - i [31] & Aspergillus fumigatus [490] & 9 & 4.3 \\
LF - 1 - i [31] & IFN-beta [1452] & 8 & 3.8 \\
LF - 2 - c [44] & Aspen [1136] & 15 & 8.4 \\
LF - 2 - c [44] & Calcium Fluoratum [1151] & 13 & 7.3 \\
LF - 2 - c [44] & Bombesina [1076] & 9 & 5.1 \\
LF - 2 - c [44] & Calcium Phosphoricum [990] & 7 & 3.9 \\
LF - 2 - e [34] & Candida albicans [556] & 16 & 6.3 \\
LF - 2 - e [34] & Aspergillus fumigatus [490] & 10 & 3.9 \\
LF - 2 - e [34] & C-RAS [1163] & 9 & 3.5 \\
LF - 2 - i [33] & Glicina [2] & 6 & 5.8 \\
LF - 2 - i [33] & Calcium Phosphoricum [990] & 5 & 4.9 \\
LF - 2 - i [33] & 25toxoplãâ sma p [2487] & 4 & 3.9 \\
LF - 2 - i [33] & AGO 1 [686] & 4 & 3.9 \\
LF - 2 - i [33] & Bombesina [1076] & 4 & 3.9 \\
LF - 2 - i [33] & H2O2 [1357] & 4 & 3.9 \\
LF - 3 - c [36] & Calcium Fluoratum [1151] & 17 & 7.9 \\
LF - 3 - c [36] & Bombesina [1076] & 14 & 6.5 \\
LF - 3 - c [36] & AGO 2 [251] & 10 & 4.6 \\
LF - 3 - c [36] & Aspen [1136] & 9 & 4.2 \\
LF - 3 - c [36] & Glutammato [2251] & 9 & 4.2 \\
LF - 3 - e [37] & Calcium Fluoratum [1151] & 17 & 8.1 \\
LF - 3 - e [37] & AGO 1 [686] & 13 & 6.2 \\
LF - 3 - e [37] & Acetilcolina [1166] & 10 & 4.8 \\
LF - 3 - e [37] & AGO 2 [251] & 9 & 4.3 \\
LF - 3 - i [35] & Calcium Fluoratum [1151] & 22 & 7.0 \\
LF - 3 - i [35] & AGO 2 [251] & 16 & 5.1 \\
LF - 3 - i [35] & Acetilcolina [1166] & 13 & 4.2 \\
LF - 3 - i [35] & GABA [1127] & 12 & 3.8 \\
LF - 3 - i [35] & Bombesina [1076] & 11 & 3.5 \\
LF - 4 - c [39] & Aspen [1136] & 15 & 8.9 \\
LF - 4 - c [39] & Acetilcolina [1166] & 11 & 6.5 \\
LF - 4 - c [39] & Calcium Fluoratum [1151] & 9 & 5.3 \\
LF - 4 - e [40] & Aspergillus niger [1267] & 6 & 4.6 \\
LF - 4 - e [40] & Candida albicans [556] & 6 & 4.6 \\
LF - 4 - e [40] & IFN-alfa [2258] & 6 & 4.6 \\
LF - 4 - i [38] & Acetilcolina [1166] & 10 & 5.2 \\
LF - 4 - i [38] & AGO 5 [1073] & 9 & 4.7 \\
LF - 4 - i [38] & Aspen [1136] & 9 & 4.7 \\
LF - 4 - i [38] & AGO 2 [251] & 8 & 4.2 \\
LF - 5 - c [45] & Calcium Fluoratum [1151] & 10 & 8.7 \\
LF - 5 - c [45] & Acetilcolina [1166] & 8 & 7.0 \\
LF - 5 - c [45] & AGO 1 [686] & 7 & 6.1 \\
LF - 5 - c [45] & Aspen [1136] & 5 & 4.3 \\
LF - 5 - c [45] & Leptina [1317] & 5 & 4.3 \\
LF - 5 - e [42] & COX-2 [679] & 7 & 4.9 \\
LF - 5 - e [42] & Candida albicans [556] & 7 & 4.9 \\
LF - 5 - e [42] & 8-OHdG [526] & 6 & 4.2 \\
LF - 5 - e [42] & C-RAS [1163] & 6 & 4.2 \\
LF - 5 - i [41] & Candida albicans [556] & 10 & 5.2 \\
LF - 5 - i [41] & 8-OHdG [526] & 8 & 4.1 \\
LF - 5 - i [41] & AGO 5 [1073] & 8 & 4.1 \\
LF - 5 - i [41] & Aspergillus niger [1267] & 8 & 4.1 \\
LF - 5 - i [41] & Aspergillus fumigatus [490] & 7 & 3.6 \\
LH - 1 - c [2] & ACTH [893] & 35 & 9.1 \\
LH - 1 - c [2] & Acetilcolina [1166] & 22 & 5.7 \\
LH - 1 - c [2] & Calcium Fluoratum [1151] & 19 & 4.9 \\
LH - 1 - c [2] & AGO 5 [1073] & 15 & 3.9 \\
LH - 1 - c [2] & AGO 2 [251] & 14 & 3.6 \\
LH - 1 - e [3] & Calcium Fluoratum [1151] & 33 & 8.8 \\
LH - 1 - e [3] & Acetilcolina [1166] & 22 & 5.9 \\
LH - 1 - e [3] & AGO 2 [251] & 15 & 4.0 \\
LH - 1 - e [3] & 3-Nitro Tirosina [361] & 14 & 3.7 \\
LH - 1 - i [1] & 8-OHdG [526] & 8 & 5.0 \\
LH - 1 - i [1] & AGO 5 [1073] & 8 & 5.0 \\
LH - 1 - i [1] & COX-2 [679] & 8 & 5.0 \\
LH - 1 - i [1] & IFN-gamma [2264] & 7 & 4.4 \\
LH - 1 - i [1] & Acido lattico [1300] & 6 & 3.8 \\
LH - 2 - c [5] & ACTH [893] & 48 & 8.2 \\
LH - 2 - c [5] & Calcium Fluoratum [1151] & 31 & 5.3 \\
LH - 2 - c [5] & AGO 5 [1073] & 30 & 5.1 \\
LH - 2 - c [5] & Acetilcolina [1166] & 25 & 4.3 \\
LH - 2 - e [6] & Acetilcolina [1166] & 18 & 6.4 \\
LH - 2 - e [6] & Calcium Fluoratum [1151] & 14 & 4.9 \\
LH - 2 - e [6] & AGO 1 [686] & 13 & 4.6 \\
LH - 2 - e [6] & Cortisolo [1056] & 11 & 3.9 \\
LH - 2 - e [6] & Rock Rose [435] & 11 & 3.9 \\
LH - 2 - e [6] & ACTH [893] & 10 & 3.5 \\
LH - 2 - e [6] & Glutammato [2251] & 10 & 3.5 \\
LH - 2 - i [4] & Candida albicans [556] & 36 & 6.7 \\
LH - 2 - i [4] & Aspergillus niger [1267] & 27 & 5.1 \\
LH - 2 - i [4] & COX-2 [679] & 20 & 3.7 \\
LH - 3 - c [8] & Acetilcolina [1166] & 31 & 11.8 \\
LH - 3 - c [8] & Calcium Fluoratum [1151] & 22 & 8.4 \\
LH - 3 - c [8] & AGO 2 [251] & 13 & 4.9 \\
LH - 3 - c [8] & Aspen [1136] & 11 & 4.2 \\
LH - 3 - c [8] & Bombesina [1076] & 11 & 4.2 \\
LH - 3 - e [9] & Acetilcolina [1166] & 51 & 9.7 \\
LH - 3 - e [9] & Calcium Fluoratum [1151] & 31 & 5.9 \\
LH - 3 - i [7] & Acetilcolina [1166] & 26 & 10.8 \\
LH - 3 - i [7] & Calcium Fluoratum [1151] & 21 & 8.7 \\
LH - 3 - i [7] & AGO 2 [251] & 10 & 4.1 \\
LH - 4 - c [11] & Acetilcolina [1166] & 43 & 11.5 \\
LH - 4 - c [11] & AGO 2 [251] & 24 & 6.4 \\
LH - 4 - c [11] & Calcium Fluoratum [1151] & 22 & 5.9 \\
LH - 4 - c [11] & Aspen [1136] & 18 & 4.8 \\
LH - 4 - c [11] & Bombesina [1076] & 18 & 4.8 \\
LH - 4 - e [12] & Acetilcolina [1166] & 24 & 9.9 \\
LH - 4 - i [10] & Calcium Fluoratum [1151] & 39 & 8.7 \\
LH - 4 - i [10] & Acetilcolina [1166] & 28 & 6.3 \\
LH - 4 - i [10] & Insulina [2282] & 20 & 4.5 \\
LH - 4 - i [10] & 3-Nitro Tirosina [361] & 19 & 4.3 \\
LH - 4 - i [10] & AGO 2 [251] & 18 & 4.0 \\
LH - 4 - i [10] & Pregnenolone [1490] & 16 & 3.6 \\
LH - 5 - c [15] & Calcium Fluoratum [1151] & 27 & 7.2 \\
LH - 5 - c [15] & AGO 2 [251] & 26 & 6.9 \\
LH - 5 - c [15] & Bombesina [1076] & 26 & 6.9 \\
LH - 5 - c [15] & Acetilcolina [1166] & 23 & 6.1 \\
LH - 5 - c [15] & Aspen [1136] & 19 & 5.1 \\
LH - 5 - e [14] & Candida albicans [556] & 33 & 5.1 \\
LH - 5 - e [14] & Aspergillus niger [1267] & 31 & 4.8 \\
LH - 5 - e [14] & Aspergillus fumigatus [490] & 25 & 3.8 \\
LH - 5 - i [13] & AGO 5 [1073] & 29 & 8.3 \\
LH - 5 - i [13] & 3-Nitro Tirosina [361] & 14 & 4.0 \\
RF - 1 - c [58] & Acetilcolina [1166] & 11 & 13.4 \\
RF - 1 - c [58] & Calcium Fluoratum [1151] & 8 & 9.8 \\
RF - 1 - c [58] & Calcium Phosphoricum [990] & 6 & 7.3 \\
RF - 1 - c [58] & Adrenalina [574] & 3 & 3.7 \\
RF - 1 - c [58] & PEtn [2185] & 3 & 3.7 \\
RF - 1 - c [58] & Profito VS - OHI [908] & 3 & 3.7 \\
RF - 1 - c [58] & Psicophyt  2 - Biogroup [210] & 3 & 3.7 \\
RF - 1 - e [47] & Aspergillus fumigatus [490] & 18 & 6.7 \\
RF - 1 - e [47] & AGO 5 [1073] & 14 & 5.2 \\
RF - 1 - e [47] & 23tãâ¨nia sãâ²lium v [2488] & 11 & 4.1 \\
RF - 1 - e [47] & Candida albicans [556] & 11 & 4.1 \\
RF - 1 - i [46] & IFN-alfa [2258] & 7 & 3.9 \\
RF - 2 - c [59] & Calcium Fluoratum [1151] & 6 & 9.0 \\
RF - 2 - c [59] & Glutammato [2251] & 6 & 9.0 \\
RF - 2 - c [59] & Acetilcolina [1166] & 5 & 7.5 \\
RF - 2 - c [59] & Dopamina [2259] & 4 & 6.0 \\
RF - 2 - c [59] & Aspen [1136] & 3 & 4.5 \\
RF - 2 - c [59] & Bombesina [1076] & 3 & 4.5 \\
RF - 2 - c [59] & Glutatione-S-Transferasi [723] & 3 & 4.5 \\
RF - 2 - e [49] & Aspergillus fumigatus [490] & 10 & 4.9 \\
RF - 2 - e [49] & Candida albicans [556] & 10 & 4.9 \\
RF - 2 - e [49] & 3-Nitro Tirosina [361] & 8 & 3.9 \\
RF - 2 - e [49] & Aspergillus niger [1267] & 8 & 3.9 \\
RF - 2 - i [48] & 3-Nitro Tirosina [361] & 7 & 9.2 \\
RF - 2 - i [48] & Acetilcolina [1166] & 5 & 6.6 \\
RF - 2 - i [48] & H2O2 [1357] & 4 & 5.3 \\
RF - 2 - i [48] & 23tãâ¨nia sãâ²lium v [2488] & 3 & 3.9 \\
RF - 2 - i [48] & AGO 2 [251] & 3 & 3.9 \\
RF - 2 - i [48] & C-RAS [1163] & 3 & 3.9 \\
RF - 2 - i [48] & Calcium Sulfuricum [701] & 3 & 3.9 \\
RF - 2 - i [48] & IL-2 [2202] & 3 & 3.9 \\
RF - 3 - c [51] & Calcium Fluoratum [1151] & 16 & 8.0 \\
RF - 3 - c [51] & Bombesina [1076] & 9 & 4.5 \\
RF - 3 - c [51] & AGO 5 [1073] & 8 & 4.0 \\
RF - 3 - e [52] & Acetilcolina [1166] & 8 & 4.4 \\
RF - 3 - e [52] & Aspen [1136] & 7 & 3.8 \\
RF - 3 - e [52] & COX-2 [679] & 7 & 3.8 \\
RF - 3 - i [50] & Calcium Fluoratum [1151] & 24 & 9.3 \\
RF - 3 - i [50] & AGO 2 [251] & 12 & 4.7 \\
RF - 3 - i [50] & Acetilcolina [1166] & 11 & 4.3 \\
RF - 3 - i [50] & Aspen [1136] & 9 & 3.5 \\
RF - 3 - i [50] & Bombesina [1076] & 9 & 3.5 \\
RF - 4 - c [54] & Calcium Fluoratum [1151] & 12 & 8.8 \\
RF - 4 - c [54] & Acetilcolina [1166] & 6 & 4.4 \\
RF - 4 - c [54] & Bombesina [1076] & 6 & 4.4 \\
RF - 4 - c [54] & 1-HL (1-Hexytol-Lysine) [1368] & 5 & 3.7 \\
RF - 4 - c [54] & Calcium Phosphoricum [990] & 5 & 3.7 \\
RF - 4 - e [55] & Aspergillus fumigatus [490] & 7 & 5.3 \\
RF - 4 - i [53] & Calcium Fluoratum [1151] & 18 & 9.1 \\
RF - 4 - i [53] & AGO 2 [251] & 13 & 6.6 \\
RF - 4 - i [53] & Dopamina [2259] & 8 & 4.1 \\
RF - 4 - i [53] & Rock Rose [435] & 7 & 3.6 \\
RF - 5 - c [60] & Calcium Fluoratum [1151] & 10 & 8.0 \\
RF - 5 - c [60] & Istamina [2293] & 7 & 5.6 \\
RF - 5 - c [60] & AGO 2 [251] & 5 & 4.0 \\
RF - 5 - c [60] & Acetilcolina [1166] & 5 & 4.0 \\
RF - 5 - c [60] & Leptina [1317] & 5 & 4.0 \\
RF - 5 - c [60] & Pregnenolone [1490] & 5 & 4.0 \\
RF - 5 - e [57] & Candida albicans [556] & 11 & 6.1 \\
RF - 5 - e [57] & Aspergillus fumigatus [490] & 9 & 5.0 \\
RF - 5 - e [57] & ACTH [893] & 8 & 4.5 \\
RF - 5 - e [57] & 8-OHdG [526] & 7 & 3.9 \\
RF - 5 - e [57] & COX-2 [679] & 7 & 3.9 \\
RF - 5 - i [56] & Candida albicans [556] & 11 & 9.2 \\
RF - 5 - i [56] & Aspergillus fumigatus [490] & 8 & 6.7 \\
RF - 5 - i [56] & 3-Nitro Tirosina [361] & 6 & 5.0 \\
RF - 5 - i [56] & Aspergillus niger [1267] & 5 & 4.2 \\
RH - 1 - c [17] & ACTH [893] & 37 & 9.9 \\
RH - 1 - c [17] & Acetilcolina [1166] & 18 & 4.8 \\
RH - 1 - c [17] & Calcium Fluoratum [1151] & 17 & 4.5 \\
RH - 1 - e [18] & Acetilcolina [1166] & 26 & 8.2 \\
RH - 1 - e [18] & Calcium Fluoratum [1151] & 26 & 8.2 \\
RH - 1 - e [18] & 3-Nitro Tirosina [361] & 14 & 4.4 \\
RH - 1 - e [18] & AGO 2 [251] & 12 & 3.8 \\
RH - 1 - e [18] & Aspen [1136] & 12 & 3.8 \\
RH - 1 - i [16] & 3-Nitro Tirosina [361] & 9 & 5.8 \\
RH - 1 - i [16] & AGO 5 [1073] & 9 & 5.8 \\
RH - 1 - i [16] & COX-2 [679] & 8 & 5.2 \\
RH - 1 - i [16] & PChoCl [2404] & 7 & 4.5 \\
RH - 1 - i [16] & Aspergillus fumigatus [490] & 6 & 3.9 \\
RH - 1 - i [16] & C-MYC [1487] & 6 & 3.9 \\
RH - 2 - c [20] & ACTH [893] & 51 & 9.9 \\
RH - 2 - c [20] & Beta estradiolo [2272] & 21 & 4.1 \\
RH - 2 - c [20] & AGO 5 [1073] & 19 & 3.7 \\
RH - 2 - c [20] & Acetilcolina [1166] & 18 & 3.5 \\
RH - 2 - e [21] & Acetilcolina [1166] & 15 & 7.9 \\
RH - 2 - e [21] & Calcium Fluoratum [1151] & 10 & 5.3 \\
RH - 2 - e [21] & Adrenalina [574] & 9 & 4.7 \\
RH - 2 - e [21] & Aspen [1136] & 8 & 4.2 \\
RH - 2 - e [21] & AGO 2 [251] & 7 & 3.7 \\
RH - 2 - e [21] & Bombesina [1076] & 7 & 3.7 \\
RH - 2 - e [21] & Dopamina [2259] & 7 & 3.7 \\
RH - 2 - i [19] & Candida albicans [556] & 32 & 7.0 \\
RH - 2 - i [19] & Aspergillus fumigatus [490] & 28 & 6.1 \\
RH - 2 - i [19] & Aspergillus niger [1267] & 16 & 3.5 \\
RH - 2 - i [19] & IFN-gamma [2264] & 16 & 3.5 \\
RH - 3 - c [23] & Acetilcolina [1166] & 19 & 10.2 \\
RH - 3 - c [23] & AGO 2 [251] & 10 & 5.3 \\
RH - 3 - c [23] & Bombesina [1076] & 8 & 4.3 \\
RH - 3 - c [23] & 3-Nitro Tirosina [361] & 7 & 3.7 \\
RH - 3 - c [23] & C-RAS [1163] & 7 & 3.7 \\
RH - 3 - e [24] & Acetilcolina [1166] & 26 & 8.8 \\
RH - 3 - e [24] & Calcium Fluoratum [1151] & 18 & 6.1 \\
RH - 3 - e [24] & AGO 2 [251] & 16 & 5.4 \\
RH - 3 - e [24] & 3-Nitro Tirosina [361] & 11 & 3.7 \\
RH - 3 - i [22] & Acetilcolina [1166] & 20 & 10.4 \\
RH - 3 - i [22] & Calcium Fluoratum [1151] & 17 & 8.8 \\
RH - 3 - i [22] & AGO 2 [251] & 16 & 8.3 \\
RH - 3 - i [22] & Bombesina [1076] & 10 & 5.2 \\
RH - 3 - i [22] & AGO 1 [686] & 8 & 4.1 \\
RH - 3 - i [22] & GABA [1127] & 7 & 3.6 \\
RH - 4 - c [26] & Calcium Fluoratum [1151] & 29 & 8.1 \\
RH - 4 - c [26] & AGO 1 [686] & 20 & 5.6 \\
RH - 4 - c [26] & Bombesina [1076] & 19 & 5.3 \\
RH - 4 - c [26] & Acetilcolina [1166] & 17 & 4.7 \\
RH - 4 - c [26] & Aspen [1136] & 15 & 4.2 \\
RH - 4 - e [27] & Acetilcolina [1166] & 29 & 13.3 \\
RH - 4 - e [27] & Calcium Fluoratum [1151] & 21 & 9.6 \\
RH - 4 - e [27] & 63 [2097] & 10 & 4.6 \\
RH - 4 - e [27] & Adrenalina [574] & 8 & 3.7 \\
RH - 4 - e [27] & Aspen [1136] & 8 & 3.7 \\
RH - 4 - e [27] & Bombesina [1076] & 8 & 3.7 \\
RH - 4 - i [25] & Acetilcolina [1166] & 36 & 9.2 \\
RH - 4 - i [25] & Calcium Fluoratum [1151] & 25 & 6.4 \\
RH - 4 - i [25] & Aspen [1136] & 16 & 4.1 \\
RH - 5 - c [30] & Acetilcolina [1166] & 22 & 6.9 \\
RH - 5 - c [30] & Calcium Fluoratum [1151] & 17 & 5.4 \\
RH - 5 - c [30] & AGO 2 [251] & 16 & 5.0 \\
RH - 5 - c [30] & Bombesina [1076] & 14 & 4.4 \\
RH - 5 - c [30] & ACTH [893] & 12 & 3.8 \\
RH - 5 - e [29] & Candida albicans [556] & 39 & 7.5 \\
RH - 5 - e [29] & Aspergillus fumigatus [490] & 22 & 4.2 \\
RH - 5 - i [28] & AGO 5 [1073] & 16 & 4.9 \\
RH - 5 - i [28] & Candida albicans [556] & 16 & 4.9 \\
RH - 5 - i [28] & Angiotensina II [561] & 13 & 4.0 \\
Thyroid - LH [61] & Acetilcolina [1166] & 2 & 9.1 \\
Thyroid - LH [61] & Adrenalina [574] & 2 & 9.1 \\
Thyroid - LH [61] & Aspen [1136] & 2 & 9.1 \\
Thyroid - LH [61] & Beta estradiolo [2272] & 2 & 9.1 \\
Thyroid - LH [61] & COX-2 [679] & 2 & 9.1 \\
Thyroid - LH [61] & PChoCl [2404] & 2 & 9.1 \\
Thyroid - LH [61] & 5-Hydroxy tryptophol [1197] & 1 & 4.5 \\
Thyroid - LH [61] & Aldosterone [2220] & 1 & 4.5 \\
Thyroid - LH [61] & IFN-alfa [2258] & 1 & 4.5 \\
Thyroid - LH [61] & IGF-1 [1415] & 1 & 4.5 \\
Thyroid - LH [61] & IL-2 [2202] & 1 & 4.5 \\
Thyroid - LH [61] & Marker 10 [2155] & 1 & 4.5 \\
Thyroid - LH [61] & SingoloUno [1650] & 1 & 4.5 \\
Thyroid - LH [61] & Sommatoria [1954] & 1 & 4.5 \\
Thyroid - LH [61] & n08 [1686] & 1 & 4.5 \\
Thyroid - LH [61] & spikeslv [2890] & 1 & 4.5 \\
Thyroid - RH [62] & PChoCl [2404] & 3 & 30.0 \\
Thyroid - RH [62] & ACo [2334] & 1 & 10.0 \\
Thyroid - RH [62] & ACoAPS [2096] & 1 & 10.0 \\
Thyroid - RH [62] & BEs [2313] & 1 & 10.0 \\
Thyroid - RH [62] & Base & 1 & 10.0 \\
Thyroid - RH [62] & Cadmio [751] & 1 & 10.0 \\
Thyroid - RH [62] & PEtn [2185] & 1 & 10.0 \\
Thyroid - RH [62] & SingoloUno [1650] & 1 & 10.0 \\
\end{longtable}
\end{footnotesize}

\subsection{Descent Shape Medians (Positive Visits Only)}

All values in this table are medians computed \emph{exclusively over positive visits} ---
visits in which an excitation anchor was detected for that point. Points with no positive
visits receive a dash (---) for all descriptor columns.
Medians are computed on the raw per-visit values; no smoothing or outlier removal is
applied.

\medskip
\noindent\textbf{How to read and interpret the columns:}

\begin{itemize}

  \item \textbf{Med.\ Slope} (\%\,step$^{-1}$): median OLS linear regression slope across
        the descent window (Module~M07). A negative value indicates a consistent downward
        trend in the post-excitation signal; a slope near zero indicates a flat or irregular
        window. The slope is the rate-of-change estimator that minimises the sum of squared
        vertical residuals between the fitted line and the window values. Undefined for
        single-marker windows (---).

  \item \textbf{Med.\ Depth} (\%): median total signal drop from the anchor value to the
        minimum value in the descent window,
        $v_{\mathrm{anchor}} - \min_{\mathrm{window}}$ (Module~M08). Larger depth indicates
        a more pronounced excitation response; a depth near zero means the signal barely
        decreases after the anchor.

  \item \textbf{Med.\ Length} (marker steps): median number of markers in the descent
        window, including the anchor (Module~M08). A longer window indicates a sustained
        post-excitation decline before a natural rise exceeding
        $\varepsilon = 5.0\,\%$ is observed or before the
        $n_{\max} = 10$ cap is reached. A length of~1 means the signal rose
        immediately at the next marker after the anchor.

  \item \textbf{Med.\ Step Decr.} (\%\,step$^{-1}$): median mean per-step decrease,
        computed as $\mathrm{depth}\,/\,(\mathrm{length} - 1)$ (Module~M08). This metric
        normalises the total drop by the window duration, providing a rate-of-descent
        estimate that is comparable across windows of different lengths. Undefined for
        single-marker windows (---).

\end{itemize}

\medskip
\noindent\textbf{Interpretation guidance.}
Anatomical points with high positive prevalence (Section~4.1), large median depth, and a
well-defined negative slope show a strong and consistent post-excitation descent across
visits.
Points with low prevalence but large depth when positive are intermittently reactive.
A high median length combined with a moderate slope may indicate a gradual, sustained
descent rather than a sharp drop.
The variance ratio (available in the detail CSV, column \texttt{variance\_ratio}) complements
these descriptors by quantifying whether signal dispersion increases or decreases after the
excitation onset.

\begin{small}
\begin{longtable}{p{4.5cm}rrrr}
\toprule
\textbf{Anatomical Point} &
\textbf{Med.\ Slope} &
\textbf{Med.\ Depth} &
\textbf{Med.\ Length} &
\textbf{Med.\ Step Decr.} \\
\midrule
\endfirsthead
\toprule
\textbf{Anatomical Point} &
\textbf{Med.\ Slope} &
\textbf{Med.\ Depth} &
\textbf{Med.\ Length} &
\textbf{Med.\ Step Decr.} \\
\midrule
\endhead
\midrule
\multicolumn{5}{r}{\small\emph{Continued on next page\ldots}} \\
\endfoot
\bottomrule
\endlastfoot
LF - 1 - c [43] & -207.50 & 185.71 & 2.0 & 98.77 \\
LF - 1 - e [32] & -171.43 & 237.50 & 2.0 & 126.92 \\
LF - 1 - i [31] & -161.20 & 208.57 & 2.0 & 114.29 \\
LF - 2 - c [44] & -147.39 & 187.12 & 2.0 & 114.44 \\
LF - 2 - e [34] & -151.92 & 237.39 & 2.0 & 134.39 \\
LF - 2 - i [33] & -160.00 & 236.36 & 2.0 & 132.50 \\
LF - 3 - c [36] & -173.50 & 231.25 & 2.0 & 143.02 \\
LF - 3 - e [37] & -142.86 & 166.67 & 2.0 & 109.09 \\
LF - 3 - i [35] & -173.16 & 257.14 & 2.0 & 141.67 \\
LF - 4 - c [39] & -141.94 & 192.98 & 2.0 & 125.00 \\
LF - 4 - e [40] & -148.08 & 230.00 & 2.0 & 116.00 \\
LF - 4 - i [38] & -171.01 & 203.57 & 2.0 & 125.00 \\
LF - 5 - c [45] & -166.67 & 178.57 & 2.0 & 118.89 \\
LF - 5 - e [42] & -162.50 & 213.35 & 2.0 & 135.16 \\
LF - 5 - i [41] & -146.67 & 202.50 & 2.0 & 115.38 \\
LH - 1 - c [2] & -183.01 & 228.34 & 2.0 & 137.09 \\
LH - 1 - e [3] & -158.62 & 194.12 & 2.0 & 120.29 \\
LH - 1 - i [1] & -160.00 & 205.88 & 2.0 & 112.50 \\
LH - 2 - c [5] & -169.23 & 211.11 & 2.0 & 133.33 \\
LH - 2 - e [6] & -175.00 & 200.00 & 2.0 & 125.00 \\
LH - 2 - i [4] & -149.17 & 195.78 & 2.0 & 115.38 \\
LH - 3 - c [8] & -178.57 & 180.00 & 2.0 & 113.33 \\
LH - 3 - e [9] & -166.09 & 192.58 & 2.0 & 116.90 \\
LH - 3 - i [7] & -152.63 & 164.71 & 2.0 & 105.56 \\
LH - 4 - c [11] & -139.57 & 173.33 & 2.0 & 104.35 \\
LH - 4 - e [12] & -160.66 & 188.46 & 2.0 & 125.32 \\
LH - 4 - i [10] & -162.50 & 177.12 & 2.0 & 114.86 \\
LH - 5 - c [15] & -153.06 & 233.33 & 2.0 & 133.33 \\
LH - 5 - e [14] & -150.00 & 212.92 & 2.0 & 115.80 \\
LH - 5 - i [13] & -166.00 & 217.86 & 2.0 & 127.78 \\
RF - 1 - c [58] & -189.54 & 188.77 & 2.0 & 90.22 \\
RF - 1 - e [47] & -173.37 & 215.38 & 2.0 & 143.21 \\
RF - 1 - i [46] & -150.57 & 234.78 & 2.0 & 108.33 \\
RF - 2 - c [59] & -137.50 & 162.07 & 2.0 & 116.24 \\
RF - 2 - e [49] & -176.47 & 260.00 & 2.0 & 156.25 \\
RF - 2 - i [48] & -188.24 & 233.33 & 2.0 & 159.77 \\
RF - 3 - c [51] & -141.67 & 206.67 & 2.0 & 113.04 \\
RF - 3 - e [52] & -180.00 & 218.75 & 2.0 & 137.25 \\
RF - 3 - i [50] & -168.42 & 204.35 & 2.0 & 136.36 \\
RF - 4 - c [54] & -163.64 & 200.00 & 2.0 & 133.75 \\
RF - 4 - e [55] & -145.65 & 226.67 & 2.0 & 115.79 \\
RF - 4 - i [53] & -156.25 & 200.00 & 2.0 & 125.73 \\
RF - 5 - c [60] & -150.00 & 235.71 & 2.0 & 117.31 \\
RF - 5 - e [57] & -140.55 & 200.00 & 2.0 & 109.52 \\
RF - 5 - i [56] & -168.75 & 187.50 & 2.0 & 142.86 \\
RH - 1 - c [17] & -176.67 & 246.67 & 2.0 & 134.21 \\
RH - 1 - e [18] & -166.29 & 188.00 & 2.0 & 120.00 \\
RH - 1 - i [16] & -163.64 & 200.00 & 2.0 & 132.74 \\
RH - 2 - c [20] & -171.13 & 231.03 & 2.0 & 153.24 \\
RH - 2 - e [21] & -166.67 & 198.33 & 2.0 & 141.18 \\
RH - 2 - i [19] & -152.38 & 211.11 & 2.0 & 127.78 \\
RH - 3 - c [23] & -160.87 & 190.48 & 2.0 & 135.00 \\
RH - 3 - e [24] & -167.50 & 216.23 & 2.0 & 127.22 \\
RH - 3 - i [22] & -166.67 & 217.39 & 2.0 & 115.00 \\
RH - 4 - c [26] & -183.33 & 183.33 & 2.0 & 117.42 \\
RH - 4 - e [27] & -151.67 & 155.04 & 2.0 & 106.97 \\
RH - 4 - i [25] & -170.00 & 200.00 & 2.0 & 120.83 \\
RH - 5 - c [30] & -161.54 & 175.00 & 2.0 & 104.00 \\
RH - 5 - e [29] & -165.38 & 226.00 & 2.0 & 141.67 \\
RH - 5 - i [28] & -162.50 & 221.05 & 2.0 & 145.61 \\
Thyroid - LH [61] & -199.00 & 54.23 & 2.0 & 54.23 \\
Thyroid - RH [62] & -180.00 & 10.00 & 1.5 & 10.00 \\
\end{longtable}
\end{small}

\end{document}
